% Author: Inigo Dela Vega
\documentclass[11pt]{article}
\usepackage[margin=1in]{geometry}
\usepackage[T1]{fontenc}
\usepackage{tabularx}
\usepackage{enumitem}
\usepackage{xcolor}
\usepackage[hidelinks]{hyperref}
\usepackage{listings}

\lstset{
  basicstyle=\ttfamily\small,
  breaklines=true,
  frame=single,
  columns=fullflexible
}

\title{Broadcasting the Bambu Lab P1S Internal Camera\\\large Research Notes and Implementation Options}
\author{Inigo Dela Vega}
\date{September 2026}

\begin{document}
\maketitle

\section{Goal}
Take the live feed from the P1S chamber camera and make it viewable somewhere other than
Bambu Studio or Bambu Handy: a browser, VLC, OBS, Home Assistant, or a remote device.

\section{Why the P1S Is Harder Than Other Bambu Printers}
Bambu printers fall into two camera groups.

\begin{center}
\begin{tabularx}{\linewidth}{@{}lXl@{}}
\hline
\textbf{Models} & \textbf{Camera interface} & \textbf{Direct RTSP?} \\
\hline
X1, X1C, X1E, P2S, H2 series & RTSPS on port 322 & Yes \\
A1, A1 Mini, P1P, P1S & Proprietary JPEG frame protocol on port 6000 & No \\
\hline
\end{tabularx}
\end{center}

Consequences for the P1S:
\begin{itemize}[itemsep=2pt]
  \item No standard RTSP URL exists, so go2rtc, VLC, and OBS cannot connect to it directly.
  \item The feed is a sequence of JPEG frames, typically around 1 to 2 fps. Expect a slideshow, not smooth video.
  \item Bambu Studio's ``Go Live'' OBS feature is not available for the P1 series.
  \item A translation layer is required: something that speaks the port 6000 protocol and re-serves the frames as MJPEG or RTSP.
\end{itemize}

\section{Approaches Others Have Used}

\subsection{BambuP1Streamer + go2rtc}
The most common community approach. BambuP1Streamer (usually run in Docker) uses Bambu's
\texttt{libBambuSource.so} to pull frames from the printer and hands them to go2rtc, which
serves a standard MJPEG stream.

\begin{lstlisting}[language={}]
streams:
  p1s: "exec:./BambuP1Streamer ./libBambuSource.so ${PRINTER_ADDRESS} ${PRINTER_ACCESS_CODE}"
\end{lstlisting}

The resulting MJPEG URL takes the form:
\begin{lstlisting}[language={}]
http://<host>:1984/api/stream.mjpeg?src=p1s
\end{lstlisting}

\textbf{Known issue:} chaining a second go2rtc instance to transcode that MJPEG into
WebRTC or H.264 (for Frigate or Home Assistant) has been reported as unreliable, with
dropped or grey frames. MJPEG out of the first instance was stable.

\subsection{Bambuddy}
A self hosted dashboard that provides MJPEG streaming from the built in printer camera and
can also ingest external network cameras. Lowest effort option for browser viewing.

\subsection{Python library: \texttt{bambu-lab-cloud-api}}
Provides a \texttt{JPEGFrameStream} class for A1/P1 printers that returns individual JPEG
frames given the printer IP and LAN access code. Suitable as the ingest layer for a
custom relay.

\begin{lstlisting}[language=Python]
from bambulab import JPEGFrameStream

with JPEGFrameStream("192.168.1.100", "access_code") as stream:
    frame = stream.get_frame()
\end{lstlisting}

\subsection{LAN Only Mode}
Some users report a noticeably faster feed when the printer runs in LAN only mode. Results
vary, and LAN only mode disables cloud features such as remote printing through Handy.

\subsection{External RTSP Camera}
For real broadcasting (Twitch, Discord, smooth remote viewing), many owners mount an
inexpensive RTSP camera instead of relying on the internal one. This bypasses the protocol
problem entirely and gives higher frame rates.

\section{Recommended Architecture}
\begin{enumerate}[itemsep=2pt]
  \item \textbf{Ingest:} a small service on a Linux host on the same LAN as the printer pulls
        JPEG frames from port 6000 (BambuP1Streamer or \texttt{JPEGFrameStream}).
  \item \textbf{Serve:} re-publish the frames as MJPEG over HTTP, either through go2rtc or a
        minimal custom server.
  \item \textbf{Consume:} open the MJPEG URL in a browser, VLC, an OBS browser or media source,
        or Home Assistant.
  \item \textbf{Remote access:} reach the relay through a VPN such as Tailscale. Do not
        port forward it.
\end{enumerate}

\section{Chosen Implementation: BambuP1Streamer}
The project will use BambuP1Streamer
(\url{https://github.com/slynn1324/BambuP1Streamer}). It pulls frames through Bambu's
plugin library and serves them with go2rtc on port 1984.

\begin{itemize}[itemsep=2pt]
  \item Tested upstream only on a P1S, on Debian 12 / amd64. ARM hosts are not supported.
  \item Pins Bambu plugin \texttt{01.04.00.15} and go2rtc \texttt{v1.9.14}. Upstream moved
        from go2rtc \texttt{v1.6.2} to \texttt{v1.9.14} on 1 June 2026, and this project
        follows upstream.
  \item Licensed AGPL-3.0.
  \item Open issues report problems with newer firmware and LAN only mode, so the working
        firmware version should be recorded once the stream is up.
\end{itemize}

On a shared network, the port is published on localhost only, so the stream is viewed on
the relay host itself:

\begin{lstlisting}[language={}]
podman build -t bambu_p1_streamer .
podman run -d --name bambu_p1_streamer --restart unless-stopped \
  -p 127.0.0.1:1984:1984 --env-file .env localhost/bambu_p1_streamer
\end{lstlisting}

Stream URL: \texttt{http://127.0.0.1:1984/api/stream.mjpeg?src=p1s}

\section{Deployment Context and Status}
\textbf{Status as of 25 September 2026: on hold pending permission.}

\begin{itemize}[itemsep=2pt]
  \item \textbf{Printer:} a P1S in the makerspace, shared by students as a whole. Its WLAN
        is not configured under a user account. Because it is shared equipment, permission
        from the makerspace is needed before using its access code or streaming its camera.
  \item \textbf{Relay host:} a ThinkPad E14 Gen 2 (x86\_64) running Omarchy, an Arch based
        distribution. Upstream was tested on Debian 12, but the container itself is Debian
        based, so the host distribution should matter little. Rootless Podman 6.1.1 is
        installed and verified.
  \item \textbf{Network:} the makerspace Wi-Fi is a large shared network with many other
        clients, unlike the home LAN assumed earlier. It may block traffic between clients,
        so reachability of the printer from the relay host must be checked first.
  \item \textbf{Done so far:} upstream cloned, both build downloads (Bambu plugin and go2rtc
        binary) confirmed available, and a private secrets file prepared outside source
        control.
  \item \textbf{Next, once approved:} enter the printer's IP and access code,
        check reachability, build, and run bound to localhost.
\end{itemize}

\section{Security Notes}
\begin{itemize}[itemsep=2pt]
  \item The printer's LAN access code grants control of the printer, not just the camera.
        Keep it in environment variables or a secrets file, never in source control.
  \item go2rtc and similar relays have little or no built in authentication. Keep them on the
        LAN and expose them only through a VPN. On a shared network such as the makerspace
        Wi-Fi, bind to localhost only.
  \item The upstream go2rtc configuration allows API requests from any origin and logs at
        debug level. Because the stream source is built from the access code, API responses
        and logs may reveal it. This is still to be verified, and logs should be redacted
        before they are shared.
  \item Using shared makerspace equipment requires the makerspace's permission.
  \item Firmware updates have changed third party local access before. Test after every update.
\end{itemize}

\section{References}
\begin{itemize}[itemsep=2pt]
  \item Bambuddy camera documentation: \url{https://wiki.bambuddy.cool/features/camera/}
  \item go2rtc issue \#1359 (P1S MJPEG to WebRTC): \url{https://github.com/AlexxIT/go2rtc/issues/1359}
  \item Bambu Lab Cloud API library: \url{https://github.com/coelacant1/Bambu-Lab-Cloud-API}
  \item Bambu forum, H2C RTSP thread: \url{https://forum.bambulab.com/t/h2c-rtsp-for-camera/256148}
  \item Bambu forum, P1P camera thread: \url{https://forum.bambulab.com/t/cant-get-camera-to-work/11084}
  \item BambuStudio issue \#1536 (RTSP request): \url{https://github.com/bambulab/BambuStudio/issues/1536}
\end{itemize}

\end{document}
